\begin{enumerate}
    \item Сложение чисел: $\mathsf{ADD} = \{(x, y, z) | x + y = z\}$ принадлежит \lsp
    \begin{enumerate}
        \item Слагаемые написаны на разных лентах: запускаем школьный алгоритм <<в столбик>>, память нужна только для хранения переноса (один бит)
        \item Слагаемые записаны на одной ленте: нужна дополнительная память для позиции, которую мы сейчас обрабатываем (логарифмическая память), ну и опять для переноса
    \end{enumerate}

    \item Умножение чисел: $\mathsf{MULT} = \{(x, y, z) | x \cdot y = z\}$ принадлежит \lsp
    
    Школьный алгоритм в чистом виде не подходит, так как подразумевает запись квадратной таблицы
    \[ x=x_n\ldots x_0, y=y_n\ldots y_0, x\cdot y=z=z_{2n}\ldots z_0\]
    \[ z_0=x_0y_0 \]
    \[ z_1=x_0y_1+x_1y_0+\text{перенос с прошлого шага} \]
    \[\ldots\]
    \[ z_k=x_0y_k+\ldots+x_ky_0+\text{перенос с прошлого шага} \]
    \[\ldots\]

    Для хранения переноса может потребоваться логарифмическая память (так как там уже не 0/1 как в сложении, а может скопиться большое число за много шагов)

    \item Проверка правильности скобочной последовательности (1 тип скобок): в каждом префиксе открывающих скобок не меньше чем закрывающих. Храним баланс (при открывающих скобках +1, при закрывающих -1). Если он уходит в минус, то отвергаем вход. Если дошли до конца слова, то надо проверить, что он равен 0. Баланс всегда не больше длины, поэтому для его хранения нужна логарифмическая память.

    \item Проверка правильности скобочной последовательности (несколько типов скобок): тут просто баланс по каждому типу не сработает, например, ([)].

    Критерий: условие на баланс соблюдается внутри каждой пары скобок. 

    Алгоритм: стартуем от открывающей скобки, заводим счетчик скобок и счетчик баланса по скобкам данного типа. Доходим до момента, когда счетчик баланса впервые обнуляется (если слово закончилось раньше, то отвергаем). Далее нужно проверить баланс скобок между найденными (можно проверить по каждому виду в отдельности, но достаточно по всем сразу). Для реализации этого алгоритма требуется небольшое число счётчиков логарифмического размера. Формальное доказательство корректности алгоритма остаётся в качестве упражнения.
\end{enumerate}